\section{Overview}\label{ch:1}

CLAS12 analyses depend on large Monte-Carlo statitics produced with the same software, geometry, and fields
configuration as the reconstructed data.

Running simulations by hand, on the OSG or the JLab farm, involves writing the submision configuration files
that contain the generator, Monte-Carlo and reconstruction conditions, chaining the pipeline outputs,
and copying results back. Given the growing number of experiments, configurations, geometry, field, and software versions,
this is error-prone and not easily reproducible.

simGrid solves this: any authenticated collaborator can request a self-consistent Monte-Carlo workflow from a browser,
with the correct software versions, field maps, vertex handling, and background merging selected for them,
and have it run at scale on the OSG without touching HTCondor directly.
The portal is the single, reproducible path from a physics request to a DST.



\subsection{Design goals and constraints}

The system is shaped by three goals and one constraint:

\begin{itemize}
  \item \textbf{Self-service:} a user with no grid or software expertise can submit a correct workflow from a form.
  \item \textbf{Operability:} the submission pipeline is self-documenting and reproducible.
  \item \textbf{Fair sharing:} no single user can monopolize the shared allocation; waiting requests are
  ordered by a history-weighted, load-aware queue.
\end{itemize}

The constraint is that \textbf{all jobs run under one shared HTCondor owner, \texttt{gemc}.} OSG allocations
and OAuth credentials are granted to that single identity, not to individual collaborators. Fair sharing
therefore cannot be delegated to HTCondor's per-user accounting; the portal must implement it itself. This
constraint is the reason the two priority systems (Chapters 10 and 11) exist and are kept separate.



\subsection{Terminology and conventions}

\begin{itemize}
  \item \textbf{Submission / cluster / subjob / job.} A \textit{submission} is a single portal request.
  When submitted to OSG, it becomes a HTCondor \textit{cluster}. A cluster contains many \textit{subjobs}
  (HTCondor processes, \texttt{\$(Process)}). Here we will use subjob and job interchangeably: a submission
  can request to run 10,000 jobs; correspondigly, the created cluster will contain 10,000 subjobs.
  \item \textbf{scard} The \textit{scard} (steering card) is the key:value dictionary the portal produces
  from the user request and stores in the database; it drives all the downstream generation.
  \item \textbf{Type-1 vs. type-2.} A \textit{type-1} submission runs an event generator on the worker node to
  produce events; a \textit{type-2} submission consumes pre-generated LUND files from a
  \texttt{/\allowbreak{}volatile/\allowbreak{}clas12/\allowbreak{}} directory, one subjob per file. The type is
  inferred from the generator field (Chapter 4).
\end{itemize}
